% PIC: @Oktavian19
\chapter{Pendahuluan}
\label{chap:pendahuluan}

\section{Gambaran Umum EVA Ecosystem}
\textbf{EVA Ecosystem} adalah platform kesehatan mental terpadu yang menghubungkan pendampingan kecerdasan buatan (\textbf{EVA AI}), pemantauan kondisi psikologis mandiri, serta layanan telekonseling bersama psikolog berlisensi dalam satu alur kerja yang aman dan berorientasi pada privasi.

Secara arsitektural, ekosistem ini memisahkan pengalaman pengguna sesuai perannya agar setiap aktor memperoleh antarmuka yang fokus, cepat, dan aman:
\begin{enumerate}[leftmargin=1.4em]
  \item \textbf{Pendampingan Mandiri Berbasis AI}: Pengguna dapat mengekspresikan kondisi emosional melalui percakapan suara maupun teks bersama EVA AI, menulis diari reflektif, serta melakukan skrining berkala menggunakan instrumen \textbf{DASS-21} (\textit{Depression, Anxiety, and Stress Scale}).
  \item \textbf{Telekonseling Klinis Terverifikasi}: Pengguna yang membutuhkan pendampingan profesional dapat menelusuri direktori psikolog terverifikasi, memesan jadwal sesi video, serta mengelola persetujuan berbagi ringkasan data klinis melalui \textit{Consent Gate}.
  \item \textbf{Tata Kelola, Verifikasi, dan Keuangan Terpusat}: Seluruh kredensial psikolog (STR/SIPP), sertifikat keaktifan, konfigurasi tarif, paket koin, pemantauan sesi, serta rekonsiliasi keuangan QRIS dan pencairan honorarium dikelola secara terpusat melalui \textbf{Web Portal dan Admin}.
\end{enumerate}

\calloutnote{Prinsip Privasi Sejak Rancangan}{Data refleksi pribadi dan catatan klinis diperlakukan dengan pembatasan akses ketat. Psikolog hanya dapat melihat ringkasan pra-sesi apabila pengguna memberikan persetujuan eksplisit pada \textit{Consent Gate}, dan Admin hanya memantau metrik agregat serta status operasional tanpa membaca isi percakapan pribadi pasien.}

\section{Peran Pengguna dalam Ekosistem}
Ekosistem EVA melayani tiga aktor utama dengan aplikasi dan tanggung jawab yang saling melengkapi, sebagaimana dirangkum pada Tabel~\ref{tab:peran-ekosistem}.

\begin{table}[htbp]
  \centering
  \caption{Ringkasan Peran, Platform, dan Tanggung Jawab Utama di EVA Ecosystem}
  \label{tab:peran-ekosistem}
  \small
  \begin{tabularx}{\textwidth}{@{}p{3.0cm}p{3.0cm}XX@{}}
    \toprule
    \textbf{Aktor} & \textbf{Aplikasi Utama} & \textbf{Fokus Aktivitas} & \textbf{Batasan Akses \& Privasi} \\
    \midrule
    \textbf{Pengguna / Pasien} & EVA Mobile (Android / Web) & Interaksi EVA AI, diari harian, skrining DASS-21, top-up koin, booking \& video telekonseling. & Memiliki kendali penuh atas \textit{Consent Gate} dan pencadangan data terenkripsi. \\
    \addlinespace[4pt]
    \textbf{Psikolog / Konselor} & EVA Counselor (Mobile) & Pendaftaran \& unggah STR/SIPP, atur jadwal \& tarif, sesi video call, catatan klinis SOAP, tarik honorarium. & Hanya melihat ringkasan pasien yang menyetujui berbagi data; tidak dapat mengakses diari mentah tanpa izin. \\
    \addlinespace[4pt]
    \textbf{Web Admin} & Web Portal \& Admin (\texttt{/admin}) & Verifikasi psikolog, penerbitan sertifikat, monitoring sesi, manajemen koin, keuangan/rekonsiliasi, CMS FAQ \& krisis. & Tidak memiliki akses ke isi diari atau transkrip konseling pribadi; hanya melihat status operasional dan metrik agregat. \\
    \bottomrule
  \end{tabularx}
\end{table}

\section{Alur Nilai Ekosistem Secara Singkat}
Keterhubungan antaraktor di EVA Ecosystem berjalan dalam siklus berikut:
\begin{itemize}[leftmargin=1.4em]
  \item \textbf{Pra-Konsultasi}: Pengguna memanfaatkan fitur pendampingan harian di \textbf{EVA Mobile}. Jika indikator skrining menunjukkan kebutuhan dukungan lanjutan atau kondisi krisis, sistem mengarahkan pengguna ke layanan \textit{SafetyNet} atau direktori psikolog.
  \item \textbf{Konsultasi Profesional}: Psikolog yang telah lolos verifikasi di \textbf{Web Admin} menerima pemesanan sesi melalui \textbf{EVA Counselor}, melaksanakan telekonseling video, lalu melengkapi catatan klinis \textbf{SOAP}.
  \item \textbf{Ekonomi Koin dan Rekonsiliasi}: Pembayaran konsultasi menggunakan \textbf{EVA Coin} yang dibeli pengguna via QRIS. Setelah kewajiban pasca-sesi terpenuhi, honorarium dilepaskan ke saldo psikolog dan dicatat dalam modul \textbf{Keuangan \& Rekonsiliasi} di Web Admin.
\end{itemize}
